% !TEX root = ../main.tex

\section{Project Workflows}
\label{sec:project-workflows}

A larger LaTeX project is more than one source file. NeoSetzer tracks common
project dependencies such as \texttt{\textbackslash input},
\texttt{\textbackslash include}, bibliography files, local letter options, and
custom document classes so that related files remain discoverable from the
workspace.

\subsection{Root documents and Magic Comments}
\label{sec:project-workflows:magic-comments}

The root declaration at the top of this child file is intentionally literal:
\begin{lstlisting}[language={[LaTeX]TeX}, caption={A child-file build declaration.}]
% !TEX root = ../main.tex
\end{lstlisting}

Magic Comments are metadata, not shell commands. NeoSetzer accepts a small,
safe subset: supported \texttt{program} values select a known LaTeX engine,
and a relative \texttt{root} value identifies a \texttt{.tex} document to
build. The application does not execute arbitrary source comments.

\subsection{Bibliography files}
\label{sec:project-workflows:bibliography}

The citation in the main document uses the bundled \texttt{references.bib}
file. Build the project twice, or let your configured build tool manage the
passes, so BibTeX can create the bibliography. The result demonstrates a real
multi-file dependency rather than a citation only written as prose.

\subsection{Structure numbering}
\label{sec:project-workflows:numbering}

The Document Structure sidebar displays standard source-level section numbers
when they can be determined reliably. It observes starred section commands and
literal section-counter changes, but it intentionally does not pretend to run
arbitrary TeX macros or custom number-format redefinitions.

The appendix at the end of \texttt{main.tex} demonstrates the common
\texttt{\textbackslash appendix} transition. In an article, the first appendix
section appears as \texttt{A}; its child headings appear as \texttt{A.1}.

\subsection{Build diagnostics}
\label{sec:project-workflows:diagnostics}

When a build fails, select an item in the build log to jump back to the
reported source line. For a harmless experiment, add an unknown command on a
new line, build once, inspect the diagnostic, and then undo the edit. Keep
experiments in this user-owned copy rather than in a valuable production
project.
